\documentclass[11pt]{article}

\usepackage[margin=1in]{geometry}
\usepackage[T1]{fontenc}
\usepackage[utf8]{inputenc}
\usepackage{lmodern}
\usepackage{microtype}
\usepackage{amsmath,amssymb,amsthm,mathtools,booktabs,array}
\usepackage{xcolor}
\usepackage[colorlinks=true,linkcolor=blue!50!black,urlcolor=blue!50!black]{hyperref}

\newcommand{\Complex}{\mathbb{C}}
\newcommand{\hh}{\mathbf{h}}
\newcommand{\QQ}{\mathbf{Q}}
\newcommand{\Gh}{\widetilde{\mathbf{G}}}
\newcommand{\x}{\mathbf{x}}
\newcommand{\I}{\mathbf{I}}
\newcommand{\norm}[1]{\left\lVert #1 \right\rVert}
\newcommand{\abs}[1]{\left\lvert #1 \right\rvert}
\newcommand{\dd}{\mathrm{d}}

\title{Direction 5: estimating $\QQ$ from channel feedback\\
\large Theory to code report (\_codex)}
\author{}
\date{2026-04-22}

\begin{document}
\maketitle

\section*{Purpose}

This report compiles the current state of the question in two senses:
\begin{enumerate}
  \item the monograph already gives a precise mathematical factorization of
  the communication object $\hh$ and the exposure object $\QQ$
  \item the AEGIS codebase already implements the route-A pipeline in which
  both are computed from the same path-level scene description
\end{enumerate}

The central conclusion is simple. Standard channel feedback gives $\hh$.
AEGIS route A computes $\QQ$ from richer path data. The theory and the code
both say the same thing: $\hh$ alone is generally insufficient to reconstruct
$\QQ$, but approximate or stale versions of $\QQ$ are still useful if their
operator-norm error is controlled.

\section*{1. Theory summary}

The monograph defines the UE channel and the exposure operator by
\begin{align}
  h_j
  &= \sum_{n:\,j(n)=j}
  \mathbf{C}_R(\hat{\mathbf{k}}_n)^H \boldsymbol{\psi}_n\,
  e^{-ik_0\hat{\mathbf{k}}_n\cdot \mathbf{r}_{\mathrm{UE}}},
  \\
  \QQ
  &= \int_\Sigma \Gh(\mathbf{r})^H \Gh(\mathbf{r})\,\dd A,
  \qquad
  P_{\mathrm{abs}}(\x)=\x^H\QQ\x,
  \\
  \Gh_j(\mathbf{r})
  &= \sum_{n:\,j(n)=j}
  \sqrt{\frac{\sigma}{4\alpha_n}}\,
  \mathbf{F}_n(\mathbf{r})\boldsymbol{\psi}_n
  e^{-ik_0\hat{\mathbf{k}}_n\cdot \mathbf{r}}.
\end{align}

These equations already expose the structural asymmetry.

\begin{enumerate}
  \item The UE channel is formed by scalar projection through
  $\mathbf{C}_R(\hat{\mathbf{k}}_n)^H$ at one location
  $\mathbf{r}_{\mathrm{UE}}$.
  \item The exposure operator is formed by Fresnel-filtered vector fields over
  the entire body surface $\Sigma$, followed by a Gram integral.
\end{enumerate}

So even before discussing estimation, the forward maps are not of the same
type. The channel is a finite vector of scalar observables. The exposure
operator is a Hermitian PSD matrix built from a surface-distributed vector
field.

\subsection*{What this implies mathematically}

The companion note
\texttt{JSAC/direction\_5\_Q\_from\_channel\_math\_codex.tex} develops the
proofs in detail. The main outcomes are:
\begin{enumerate}
  \item $\hh$ does not generically identify $\QQ$
  \item there is no nontrivial $\hh$-only bound on
  $\rho = \hh^T\QQ\hh^*/(\norm{\hh}^2\lambda_{\max}(\QQ))$ without extra
  priors on the admissible family of operators
  \item if $\widehat{\QQ}$ approximates $\QQ$ with
  $\varepsilon_Q = \norm{\QQ-\widehat{\QQ}}_2$, then
  \begin{equation}
    \abs{\x^H\QQ\x-\x^H\widehat{\QQ}\x} \le \norm{\x}^2 \varepsilon_Q
  \end{equation}
  for every precoder $\x$
\end{enumerate}

That third item is the key bridge from theory to control. It says one does
not need exact inversion of $\QQ$ to design safe precoders. One needs a bound
on the operator error.

\section*{2. Code realization in AEGIS}

The codebase already mirrors the monograph structure quite closely. The table
below is the important map.

\begin{center}
\begin{tabular}{p{0.22\linewidth} p{0.34\linewidth} p{0.34\linewidth}}
\toprule
Object & Theory role & Code realization \\
\midrule
Propagation paths & shared latent data
$\{ \hat{\mathbf{k}}_n,\boldsymbol{\psi}_n,j(n)\}$ &
\texttt{src/aegis/paths.py} \\
UE channel $\hh$ & signal branch, scalar projection of paths &
\texttt{src/aegis/mimo/channel.py} \\
Field channel $G(\mathbf{r})$ & free-space vector field on body &
\texttt{src/aegis/coherent/field\_channel.py} \\
Body channel $\Gh(\mathbf{r})$ & Fresnel and depth filtered field &
\texttt{src/aegis/coherent/body\_channel.py} \\
Exposure operator $\QQ$ & Gram integral of $\Gh^H\Gh$ &
\texttt{src/aegis/coherent/exposure\_operator.py} \\
Alignment $\rho$ & MRT exposure metric &
\texttt{compute\_rho()} in exposure operator module \\
ECBF solver & constrained optimal precoder &
\texttt{src/aegis/coherent/ecbf.py} \\
Level 7 and 8 kernels & coherent map and ECBF map &
\texttt{src/aegis/kernels/level7\_coherent.py},
\texttt{level8\_ecbf.py} \\
\bottomrule
\end{tabular}
\end{center}

\subsection*{Where the latent variables live}

\texttt{PropagationPaths} stores exactly the shared path-level data that
matters for this question: arrival directions \texttt{k\_hat}, complex
polarisation-amplitude vectors \texttt{psi}, and originating element indices
\texttt{element\_index}. In other words, the code already reflects the
monograph's view that both signal and exposure are functions of the same
underlying path representation, but through different observation operators.

\subsection*{Where $\hh$ is computed}

The communication channel is computed in
\texttt{src/aegis/mimo/channel.py}. The function
\texttt{compute\_channel\_vector(...)} explicitly forms:
\begin{enumerate}
  \item the UE receive pattern via the half-wave dipole effective length
  \item a per-path scalar coefficient
  \[
    c_n = \overline{\mathbf{C}_R(\hat{\mathbf{k}}_n)} \cdot
    \boldsymbol{\psi}_n\,
    e^{-ik_0\hat{\mathbf{k}}_n\cdot \mathbf{r}_{\mathrm{UE}}}
  \]
  \item the final array channel $\hh = A^T c$
\end{enumerate}

This is exactly the scalar-collapse story in the monograph. The UE sees one
number per path after receive-pattern projection. The rest of the vector
structure is discarded.

\subsection*{Where $\QQ$ is computed}

The body-surface channel is built in
\texttt{src/aegis/coherent/body\_channel.py} by:
\begin{enumerate}
  \item computing Fresnel coefficients and TE/TM bases for each
  body-triangle and path pair
  \item applying the Fresnel operator to each path vector
  \item applying the depth weight $\sqrt{\sigma/(4\alpha)}$
  \item multiplying by phase factors
  \item accumulating contributions by antenna element
\end{enumerate}

Then \texttt{compute\_exposure\_operator(...)} in
\texttt{src/aegis/coherent/exposure\_operator.py} forms
\begin{equation}
  \QQ = \sum_m A_m\,\Gh_m^H \Gh_m,
\end{equation}
enforces Hermitian symmetry numerically, and provides the eigendecomposition
and $\rho$ calculation.

This is not an approximation to the theory. It is the theory, discretized on
the body mesh.

\section*{3. Where route A already exists in code}

The most relevant implementation for direction 5 is the MIMO pipeline in
\texttt{src/aegis/mimo/compute.py}. For each user, it already does the
following:
\begin{enumerate}
  \item generate center paths toward the device
  \item expand them to per-element paths
  \item compute the body channel $\Gh$
  \item compute the exposure operator $\QQ$
  \item compute the communication channel $\hh$
\end{enumerate}

That ordering matters. The code does not try to infer $\QQ$ from $\hh$. It
computes both from a richer path model. That is exactly route A in operational
form.

The same file then hands the per-user operators to the MIMO precoder stage.
So from a systems perspective, AEGIS already contains the infrastructure
needed to study:
\begin{enumerate}
  \item how sensitive $\QQ$ is to pose or geometry changes
  \item how often $\QQ$ needs to be refreshed relative to $\hh$
  \item whether robust margins based on stale $\widehat{\QQ}$ are sufficient
\end{enumerate}

\section*{4. What the code confirms about the research question}

The code confirms four points that are easy to miss if one thinks only in
communications notation.

\subsection*{Point 1. The shared object is not $\hh$, it is the path set}

Both the signal branch and the exposure branch start from
\texttt{PropagationPaths}. So the common latent variable is not the channel
vector. It is the richer per-path tuple
\[
  (\hat{\mathbf{k}}_n,\boldsymbol{\psi}_n,j(n)).
\]
This immediately explains why standard channel sounding is insufficient.

\subsection*{Point 2. The missing information is geometric and polarimetric}

The channel computation uses a single receive pattern and a single UE point.
The exposure computation uses body normals, triangle centroids, Fresnel bases,
depth coupling, and surface integration. The missing information is therefore
not a small correction term. It is most of the geometric operator.

\subsection*{Point 3. The code already supports stale-$\QQ$ experiments}

Because $\QQ$ is assembled explicitly from mesh pose and paths, one can perturb
body orientation, body position, or scene geometry and directly measure
\[
  \norm{\QQ(\xi)-\QQ(\xi_0)}_2,
  \qquad
  \abs{\lambda_{\max}(\QQ(\xi))-\lambda_{\max}(\QQ(\xi_0))},
  \qquad
  \abs{\rho(\xi)-\rho(\xi_0)}.
\]
This makes the staleness question unusually tractable.

\subsection*{Point 4. Level 8 is already the right evaluation loop}

The level-8 kernel computes $\QQ$, solves ECBF, and returns the resulting
absorbed power density map. So the effect of a stale or approximate
$\widehat{\QQ}$ can be tested by replacing the operator used in the solver and
comparing:
\begin{enumerate}
  \item feasibility under the true $\QQ$
  \item signal loss relative to fresh $\QQ$
  \item degradation in absorbed-power control
\end{enumerate}

\section*{5. The most promising next experiments}

The code and theory together suggest that the most useful next report is not
``recover $\QQ$ from $\hh$''. It is ``measure the regularity of $\QQ$ under
controlled perturbations''.

\subsection*{Experiment A. Orientation Lipschitz constant}

Fix scene, array, frequency, and user position. Rotate the body by small
angles $\theta$ around one or more axes. For each pose, compute
$\QQ(\theta)$. Estimate
\begin{equation}
  L_\theta \approx
  \max_{\theta\neq\theta'}
  \frac{\norm{\QQ(\theta)-\QQ(\theta')}_2}{\abs{\theta-\theta'}}.
\end{equation}
Then the first-order control law is
\begin{equation}
  \abs{\x^H\QQ(\theta)\x-\x^H\QQ(0)\x}
  \le \norm{\x}^2 L_\theta \abs{\theta}.
\end{equation}

\subsection*{Experiment B. Refresh interval}

If the body rotates with angular speed bounded by $\omega_{\max}$, then over
an interval $\Delta t$ the operator drift is approximately bounded by
\begin{equation}
  \norm{\QQ(t+\Delta t)-\QQ(t)}_2
  \lesssim L_\theta \omega_{\max}\Delta t.
\end{equation}
This yields a practical refresh rule for route A:
\begin{equation}
  \Delta t_{\max}
  \approx \frac{\varepsilon_Q}{L_\theta \omega_{\max}},
\end{equation}
where $\varepsilon_Q$ is the tolerated operator error.

\subsection*{Experiment C. Geometry mismatch}

Hold body pose fixed and perturb scene geometry: remove or move scatterers,
perturb reflector positions, or change the set of dominant rays. Then compare
the stale operator $\widehat{\QQ}$ against the fresh one. This directly tests
the engineering question in the markdown note: does route-A staleness create
real compliance risk or mostly performance loss.

\subsection*{Experiment D. Robust ECBF margin}

Take a stale operator $\widehat{\QQ}$ and solve ECBF against a tightened bound
\begin{equation}
  \x^H\widehat{\QQ}\x \le P_{\mathrm{abs}}^{\max} - P\varepsilon_Q.
\end{equation}
Then evaluate the resulting precoder under the true operator $\QQ$. This tests
whether a simple operator-norm margin is already enough for conservative
control.

\section*{6. What I would claim in the paper right now}

Based on both the monograph and the implemented code, I think the strongest
defensible claim is:

\begin{quote}
The exposure operator $\QQ$ is not identifiable from standard channel
feedback alone because the channel vector $\hh$ is only a scalar projection of
the underlying path-level vector field. However, AEGIS already computes both
$\hh$ and $\QQ$ from a common path representation, making route-A estimation a
practical basis for studying staleness. The relevant technical quantity is the
operator error $\norm{\QQ-\widehat{\QQ}}_2$, because it directly controls
absorbed-power error, worst-case exposure drift, and robust compliance
margins.
\end{quote}

That is a stronger and cleaner story than trying to force a channel-only
inverse problem that the current mathematics does not support.

\section*{Bottom line}

The codebase already contains the exact machinery needed for the most
interesting version of direction 5:
\begin{enumerate}
  \item route-A construction of $\hh$ and $\QQ$ from common paths
  \item explicit coherent body channel construction
  \item explicit exposure operator assembly
  \item explicit ECBF optimization
\end{enumerate}

So the next deliverable should be empirical and operator-theoretic:
measure how $\QQ$ moves, not whether $\hh$ magically contains it.

\end{document}
